% Drawing 148/01-18 -- folder 148, batch 01, archive page 20 (PDF page 21).
% Transcription: au crayon, sans texte : un tétraèdre (l'arête cachée en tirets), entouré de six traits rayonnants ; à droite, un petit polyèdre (octaèdre, ou deux tétraèdres croisés) traversé par un axe.
% Kind: pencil tetrahedron with six radiating strokes; small polyhedron with an axis. Method 2 (schematic redraw from the facsimile). Status: ok.
% Notes: Faint pencil on the back of a listing (seen with the contrast raised). Left: a tetrahedron with four vertices, five edges solid and the hidden edge dashed, all drawn between the same four points. Around it, long strokes radiate without touching it: up-left, up (a narrow doubled stroke), down-left, down-right and down; a short stroke stands at the far left; and the stroke going up-right is the axis of the second figure, which runs on beyond it. Right (figure 2): a small polyhedron of overlapping strokes, drawn stroke by stroke as it stands: a top vertex joined by a long stroke to a left vertex and by another to a bottom vertex, a curved stroke from the top to an upper right vertex and on to a right vertex, from the left vertex one stroke to a middle point and one to the bottom, from the right vertex one stroke to the bottom and one back to the middle point; crossed by an oblique axis, heavily gone over inside the figure and running on thinly to the upper right edge of the crop. Left out: the typed digits of the listing showing through at the lower left. Redrawn 2026-11-12 (issue 128). 2026-11-13, Claude Opus 5.5, issue 147, overlay check: not compiled (no TeX in the session); the TikZ was rasterised by a local interpreter and laid over the 236 ppi scan rendered at 400 dpi, contrast raised. Redrawn on the ink: every stroke fitted on its ink (least squares, or points moved onto the ink), so each lies on its ink to within about half a unit. Changed: the radiating strokes are where they are on the page (the upper left one starts much higher, the upper right one is the axis of the polyhedron); the short stroke at the far left is added; the polyhedron is drawn as its strokes stand, not as a quadrilateral with its diagonals. Still a reading: which solid the polyhedron is (octahedron or two crossed tetrahedra, as the transcription says) is not decided by its strokes; the redraw does not decide it. Disagreement with the transcription: six radiating strokes are counted only if the up-right one is counted, and it is the axis that crosses the polyhedron; there is a seventh, short, at the far left.
% Coordinates: 1 unit = 1 mm = 1/100 of the crop width, y downwards as on the crop (bbox in page fractions: [0.26, 0.13, 0.7, 0.36]).
% Redrawn by Claude (claude-fable-5-2), issue #128, 2026-11-12; not checked by a human; not compiled (no TeX in the session); overlay-checked on the scan by Claude Opus 5.5, issue #147, 2026-11-13, not compiled (no TeX in the session).
% Rights: reproduces part of the work -- behind RENDER_TRANSCRIPTS, never published by itself.
\documentclass[tikz,border=3pt]{standalone}
\usepackage{amsmath,amssymb}
\usetikzlibrary{patterns,arrows.meta}
\tikzset{pencil/.style={black!45}, struck/.style={black!35}, mpt/.style={circle,draw,inner sep=1.3pt}, pt/.style={circle,fill,inner sep=0.9pt}, opt/.style={circle,draw,fill=white,inner sep=1.1pt}, Rd/.style={red!75!black}, Gn/.style={green!45!black}, Br/.style={brown!70!black}}
\begin{document}
\begin{tikzpicture}[x=1mm,y=-1mm,line cap=round,line join=round,font=\small]
% tetrahedron: five edges solid, the hidden one dashed
\coordinate (T) at (51.7,30.1);
\coordinate (L) at (39.55,46.3);
\coordinate (R) at (62.1,41.7);
\coordinate (B) at (54.9,52.3);
\draw[pencil,thick] (T) -- (L) -- (B) -- (R) -- cycle;
\draw[pencil,thick] (T) -- (52.28,32.99) -- (53.07,37.99) -- (53.99,42.97) -- (54.76,46.95) -- (55.3,49.93) -- (B);
\draw[pencil,thick,dashed] (L) -- (R);
% strokes radiating round it (the top one doubled on the page), and a short stroke at the far left
\draw[pencil] plot[smooth, tension=0.55] coordinates {(7.94,5.86) (13.92,12.08) (20.04,17.97) (25.91,24.58) (31.92,31.07) (38.11,37.91)};
\draw[pencil] plot[smooth, tension=0.55] coordinates {(54.72,0.33) (55.31,6.02) (55.45,11.99) (55.6,18.01) (55.54,24.5)};
\draw[pencil] (31.88,49.62) -- (3.83,69.22);
\draw[pencil] (63.35,51.91) -- (87.74,73.55);
\draw[pencil] (51,59.91) -- (50.77,73.72);
\draw[pencil] (7.53,37.06) -- (11.75,33.12);
% the small polyhedron, as its strokes stand
\draw[pencil] (78.13,20.09) -- (69.3,34.83);
\draw[pencil] plot[smooth, tension=0.55] coordinates {(78,20.57) (80.44,22.08) (82.5,23.6) (84.53,25.27)};
\draw[pencil] (83.27,24.12) -- (88.75,29.04);
\draw[pencil] (77.93,20.38) -- (83.39,37.36);
\draw[pencil] (69.2,34) -- (79.8,31.24);
\draw[pencil] (69.25,33.9) -- (82.49,38.34);
\draw[pencil] (88.57,28.38) -- (81.58,38.25);
\draw[pencil] (78.15,30.71) -- (88.88,28.65);
% the axis crossing it, gone over heavily inside it, running on to the upper right
\draw[pencil] (69.55,38.84) -- (100.03,10.1);
\draw[pencil,line width=1.6pt] (74.66,36.13) -- (84.45,24.76);
\end{tikzpicture}
\end{document}
